% !TEX root = ../main.tex
\section{Industrial Context and Evaluation Scope}
\label{sec:context}
\label{sec:industrial-context}

\subsection{The Deployment Decision We Had to Make}
\label{sec:deployment-decision}

The motivation for this study is an adoption decision rather than a
leaderboard exercise.  Huawei engineering teams working on accelerator
operators, distributed training, and Kubernetes-based platform software
are evaluating whether commercial coding agents can safely take on
\emph{feature-level} work, not just code completion or local bug fixes.
The practical question is concrete: can an agent produce a first-draft
pull request close enough to a senior engineer's implementation that a
reviewer could land it with minor edits?

This deployment question decomposes into four research questions:

\begin{enumerate}
  \item \textbf{RQ1: Capability boundary.} Can current agents complete
  production-equivalent feature PRs without human intervention?
  \item \textbf{RQ2: Task-shape effect.} Does performance degrade
  gradually as feature size grows, or is there a cliff for cross-module
  work?
  \item \textbf{RQ3: Harness leverage.} How much can richer task
  specifications help without changing the underlying model or agent
  framework?
  \item \textbf{RQ4: Deployment policy.} Which task classes should be
  routed to agents today, and which should remain human-owned?
\end{enumerate}
These questions shape the benchmark design, prompt variants, and
deployment guidance reported in the rest of the paper.

\subsection{Why Generic Benchmarks Were Insufficient}
\label{sec:why-not-swe-bench}

Before building our own evaluation, we considered existing coding-agent
benchmarks. SWE-Bench~\cite{jimenez2024swebench} and its
variants~\cite{chowdhury2024swebenchverified,swebenchpro2025} evaluate
real repository patches, but the dominant workflow is bug fixing with
an executable oracle.  Feature implementation in our setting has a
different acceptance condition: the agent must discover every file that
belongs to the feature, add tests where appropriate, and follow the
project's architectural conventions. We therefore needed an evaluation
that:
\begin{itemize}
  \item mirrored the \emph{structure} of the work our engineers actually do
    (multi-file features, test additions, cross-subsystem refactors),
  \item drew tasks from \emph{our} codebases so that language, API, and
    architecture matched deployment reality, and
  \item used a verdict rule robust to single-evaluator bias.
\end{itemize}

Existing public benchmarks cover parts of this space, but none provided
all three properties simultaneously for our deployment decision.

\subsection{Codebase and Task Scope}
\label{sec:codebase-scope}

The benchmark pairs 12 tasks mined from Huawei-maintained production
repositories with 50 tasks sampled from widely-used open-source
repositories so that results generalise beyond one organisation's coding
style.

\paragraph{Huawei and platform tasks (12 tasks).}
The 12 industrial tasks span three engineering areas selected with our
engineering partners as representative consumers of an internal agent
deployment:
\begin{itemize}
  \item \textbf{CANN/\texttt{torch\_npu} accelerator stack} (\textsc{C1}--\textsc{C5}):
    performance-sensitive C++/AscendC kernels and Python integration
    changes implementing new operators or refactoring existing ones;
    typical changes touch kernel code, host-side launch wrappers,
    operator registration tables, integration code, and unit tests.
  \item \textbf{MindSpeed training framework} (\textsc{M1}--\textsc{M3}):
    Python changes to the distributed-training runtime; typical tasks
    modify scheduling, memory accounting, or loss-scaling logic and
    require passing the MindSpeed integration suite.
  \item \textbf{Kubernetes platform} (\textsc{K1}--\textsc{K4}): feature
    work on the upstream open-source Kubernetes codebase (Go), including
    API server extensions, controller changes, and generator-driven
    conversion code spanning 10s of files per PR.
\end{itemize}

\paragraph{Open-source control set (50 tasks).}
The LongBench tasks (\textsc{T01}--\textsc{T50}) are merged PRs sampled
from 15 widely-deployed open-source projects (Kafka, CPython, Airflow,
Django, gRPC, Spring, Kubernetes OSS, and others). They provide an
out-of-organization control set with similar multi-file feature shapes
but without Huawei-specific APIs.

\subsection{Task Curation Funnel}
\label{sec:task-funnel}

Raw PR archives are unsuitable as benchmark tasks because many PRs are
trivial, non-reproducible, or depend on unstated team context. We
therefore used the same four-stage curation funnel for both the Huawei
and LongBench subsets:

\begin{enumerate}
  \item \textbf{Feature-work filter.} Retain PRs that implement a
  behaviour change, feature extension, or cross-file refactor; drop
  typo fixes, dependency bumps, formatting-only changes, and
  documentation-only changes.
  \item \textbf{Reproducibility filter.} Retain tasks whose base
  repository can be checked out and sanitised without requiring
  unreleasable ground-truth history, hidden remotes, or production
  credentials.
  \item \textbf{Reviewability filter.} Retain tasks whose expected
  outcome can be compared against a merged implementation with concrete
  files, tests, or behaviour, rather than an open-ended design
  discussion.
  \item \textbf{Diversity filter.} Stratify by repository family,
  language, file count, and diff size to avoid a benchmark dominated by
  one project or one task style.
\end{enumerate}

The final benchmark contains 62 tasks: 12 Huawei/platform tasks and 50
LongBench tasks. We intentionally prefer a smaller, auditable benchmark
over a large automatically harvested set because the study's purpose is
deployment guidance, not training-data scale.

\subsection{Deployment-Equivalence Criterion}
\label{sec:deployment-criterion}

We score agent outputs under a production-equivalence criterion: a PASS
means the generated diff is close enough to the merged human PR that a
reviewer could plausibly accept it after ordinary review. This criterion
combines functional correctness, completeness of the edited scope, and
behavioural equivalence to the accepted design. It is intentionally
stricter than test-only resolution because industrial feature work often
has multiple files, generated artifacts, and reviewer conventions that
are not fully captured by a single fail-to-pass test. Section~\ref{sec:study-design}
defines the evaluator panel and precise rubric.

\subsection{Scope and Non-Goals}
\label{sec:scope}

We do not measure productivity impact, developer satisfaction, or
long-term maintenance outcomes; those require live deployment studies.
We also do not study pair-programming loops with intermediate human
feedback. Every run is fully autonomous because that is the strongest
claim a procurement or platform team would need before routing feature
work to an agent. Finally, we deliberately exclude UI/design tasks,
documentation-only PRs, and PRs that were reverted or superseded.
